problem:  
Let \(B_r^n \subset \mathbb{S}^n\) be the geodesic ball of radius \(r < \pi/2\), and consider the class of compact, smoothly embedded, free boundary minimal submanifolds \(\Sigma^k \subset B_r^n\) (with \(\partial \Sigma \subset \partial B_r^n\), meeting it orthogonally).  

Define the **deficit functional**  
\[
\Delta_r(\Sigma) = \frac{\operatorname{Vol}_k(\Sigma)}{\operatorname{Vol}_k(D_r^k)} - 1 ,
\]
where \(D_r^k = B_r^n \cap \mathbb{S}^k\) is a totally geodesic \(k\)-disk.  
Known rigidity theorems (e.g. Brendle and others) imply \(\Delta_r(\Sigma) \ge 0\), with equality only for \(D_r^k\).

**Problem (Quantitative Spherical Stability Conjecture):**  
Is there a universal dimensional constant \(C(n,k,r)>0\) such that for all compact, embedded, free boundary minimal submanifolds \(\Sigma^k \subset B_r^n\),
\[
\mathrm{dist}_{C^1}(\Sigma, D_r^k) \le C(n,k,r) \, \Delta_r(\Sigma)^{\alpha},
\]
for some optimal exponent \(\alpha > 0\), where \(\mathrm{dist}_{C^1}\) denotes the distance in the \(C^1\) topology after optimal congruence in \(\mathbb{S}^n\)?  
Equivalently, quantify **how close** a free boundary minimal submanifold with almost minimal volume must be to the totally geodesic disk.

Why is it a "good" problem:  
This problem moves beyond the known rigidity theorem to a **quantitative stability** question, seeking to measure *how strongly* the totally geodesic disk minimizes volume among free boundary minimal submanifolds.  

1. **New Territory Beyond Known Results:**  
   While the rigidity inequality itself is established, sharp quantitative stability (with explicit exponents or optimal constants) is unknown even for \(k=2\), and there are virtually no results for higher dimensions or codimensions. This is exactly the next frontier after classical qualitative rigidity.

2. **Bridge Across Fields:**  
   Resolving this question requires ideas from **geometric analysis** (second variation, stability operators, and index theory), **PDE regularity theory** (quantitative control from the minimal surface equation), and **metric geometry** (geometry of nearly minimizing submanifolds). Hence it connects differential geometry, analysis, and geometric measure theory in a unified framework.

3. **Profound Simplicity with Deep Consequences:**  
   The statement is simple—control \(C^1\)-distance by a volume deficit—but its proof would likely demand major advances in understanding the spectral gap of the Jacobi operator and the nonlinear stability of minimal submanifolds under free boundary conditions.  
   Success here could inspire analogous quantitative results in Euclidean, hyperbolic, or anisotropic settings, enriching rigidity theory.

Thus, this “Quantitative Spherical Stability Conjecture” is a natural and compelling extension of classical rigidity theorems, combining simplicity of form with profound analytical and geometric depth, making it a genuinely “good” open research problem.